\section{Introduction}
\label{sec:intro}

Estimates of the time-varying reproduction number have become a routine output of infectious disease surveillance. They were used to track the transmission of Ebola virus disease in West Africa \citep{nouvellet2018simple}, to assess non-pharmaceutical interventions against COVID-19 \citep{flaxman2020estimating} and to describe the early spread of SARS-CoV-2 \citep{kucharski2020early}. They are now produced routinely for many countries \citep{huisman2022estimation}, with practical guidance on their estimation \citep{gostic2020practical, steyn2025primer}. Their value for situational awareness depends on how precisely they can be estimated from the data that surveillance systems actually collect.

The instantaneous reproduction number, $R_t$, is the expected number of secondary infections generated by an individual infected at time $t$ if conditions remained as they were at that time \citep{fraser2007estimating}. It is usually estimated from incidence through the renewal equation, which weights past incidence by the generation interval \citep{wallinga2004different, cori2013new, thompson2019improved}. The generation interval is the time between infection of a primary case and infection of its secondary cases. For a given generation interval, the posterior uncertainty of such estimates shrinks as the number of cases in the estimation window increases \citep{cori2013new}. Parag and colleagues improved estimation at low incidence \citep{parag2021improved}. They also quantified in closed form the information that incidence carries about $R_t$, including the loss caused by time-varying reporting fractions \citep{parag2022quantifying}, and derived the limits this information places on timely interventions \citep{parag2025asymmetric}. Extensions handle temporally aggregated counts \citep{nash2023estimating, ogigittins2024simulation} and the difference between the reproduction number of infections and the apparent reproduction number of a delayed observation series \citep{eales2024differences}.

These precision results measure error against the realised value of $R_t$ in the period being estimated, and that error vanishes as case counts grow. Many uses of $R_t$ estimates refer instead to a quantity that persists across periods, such as the typical level of transmission in a season or the value expected next week. Policy evaluations likewise average over weeks with different climatic conditions, school calendars and contact patterns, all of which modulate influenza transmission \citep{cauchemez2008estimating, shaman2009absolute, shaman2010absolute}. Between-period variation in $R_t$ is shared by all chains of transmission in a population at a given time. Population ecology distinguishes such environmental stochasticity from demographic stochasticity, which arises from independent individual outcomes \citep{engen1998demographic, rohani2002determinism, lande2003stochastic}, and mechanistic epidemic models have included both forms \citep{breto2009time}. Demographic variance per individual falls with population size, whereas environmental variance does not, and the population size at which the two contributions are equal is a standard summary of extinction risk \citep{lande1993risks, lande2003stochastic}. Overdispersed offspring distributions \citep{lloydsmith2005superspreading, endo2020estimating} and incomplete reporting add to the demographic side of this partition.

Surveillance data add a third source of variation. The proportion of infections that appear in a count changes over time with testing practice, care seeking and the completeness of facility reporting \citep{azmon2014estimating, stoner2019hierarchical, parag2022fundamental}. In the United States, hospitals were required to report influenza admissions to the National Healthcare Safety Network (NHSN) from February 2022 \citep{cdc2022flusight, frutos2025evaluation}. Reporting was voluntary from May to October 2024 and required again from November 2024 \citep{cms2024reporting, cdc2026nhsn}. Hospital admission series of this kind are increasingly used for forecasting and real-time assessment \citep{mathis2024evaluation}. Variation in reporting completeness that persists within a week acts on every case in that week and therefore behaves like environmental variation, not like sampling noise.

Sample size calculations for surveillance usually assume that more cases give proportionally more precision \citep{cochran1977sampling}, and analyses of reproduction number estimation have concentrated on the informativeness of incidence for period-specific values \citep{parag2022quantifying, bhatt2023semi}. A variance that is the sum of a term shared within a unit and a term that falls with the number of observations in the unit is, however, familiar from several fields. It is the variance of a cluster mean in two-stage sampling \citep{cochran1977sampling}, and it appears in sample size calculations for cluster-randomised trials. There the between-cluster coefficient of variation sets the number of clusters needed and the cost ratio sets the optimal cluster size \citep{hayes1999simple, raudenbush1997statistical}. Designs for detecting trends in repeated ecological surveys face the same structure, because variation shared by all sites in a year limits what additional sites can achieve \citep{urquhart1999designs, kincaid2004structure}. State-space models of population dynamics show that process and observation variance are often difficult to separate from a single series \citep{knape2008estimability, augermethe2016state}. Stochastic epidemic models fitted to incidence include environmental noise in transmission alongside demographic noise \citep{breto2009time, he2010plug, dureau2013capturing, king2015avoidable}. To our knowledge, this partition has not been carried over to the precision and design of surveillance for the reproduction number. In surveillance the term that falls with case volume depends on the mean reproduction number through overdispersed offspring and binomial reporting, volatile reporting completeness adds a term that does not fall with case volume, and a budget is spread over several periods.

An empirical test also faces a measurement problem. Weekly growth statistics computed from admissions are apparent quantities whose variance depends on the delay from infection to admission and on how each statistic aggregates growth, so their intercepts cannot be read directly as transmission variance. Observations from the same calendar week share national influences, and interval estimates that treat them as independent are too narrow \citep{owen2007pigeonhole, cameron2011robust}.

Here we carry the ecological variance partition over to case-count estimates of the reproduction number under environmental variation in $R_t$, negative binomial offspring and binomial reporting, and derive the balance scale at which environmental and sampling contributions are equal. The partition concerns the mean reproduction number across periods and one-step prediction; for the period-specific value $R_t$ the sampling error still vanishes as counts grow. We examine its consequences for reporting heterogeneity and for dividing a fixed budget between cases per period and number of periods. We then fit the relative loss of two apparent weekly statistics in NHSN influenza admissions and map their intercepts to the latent environmental variance by a two-step inversion through a data-matched simulation, keeping the apparent scale of each statistic distinct from the latent scale of the model (Table~\ref{tab:notation}). Finally, we decompose the intercept into national-week and state-specific components and compare the mandatory and voluntary reporting regimes.
